\documentclass[a4paper,11pt]{article}

\usepackage[czech]{babel}
\usepackage[T1]{fontenc}
\usepackage[utf8]{inputenc}
\usepackage{amsmath,amssymb}
\usepackage{geometry}
\usepackage{enumitem}

\geometry{a4paper,margin=17mm}
\pagestyle{empty}
\setlength{\parindent}{0pt}
\setlength{\parskip}{4pt}
\setlist[enumerate]{leftmargin=*,itemsep=4pt,topsep=4pt}

\newcommand{\TestTitle}{Komplexní čísla}
\newcommand{\TestHeader}[1]{%
    \begin{center}
        {\Large\bfseries TEST -- \TestTitle}\\[2mm]
        {\bfseries Varianta #1}
    \end{center}
    \vspace{2mm}
    Jméno a příjmení: \rule{9cm}{0.4pt}\hfill Třída: \rule{2cm}{0.4pt}\\[3mm]
    Datum: \rule{3cm}{0.4pt}\hfill Počet bodů: \rule{2cm}{0.4pt}\hfill Známka: \rule{2cm}{0.4pt}\\[2mm]
    \hrule\vspace{4mm}
}
\newcommand{\AnswerSpace}[1]{\vspace{#1}}

\begin{document}
\TestHeader{A}

% ZADÁNÍ VARIANTY A
\begin{enumerate}
    \item \((2+3i)^2\)
    \AnswerSpace{17mm}
    \item \(\left|\frac{2+3i}{1-i}\right|\)
    \AnswerSpace{24mm}
    \item \(x^2-5x+5=0\), řešte v \(\mathbb C\).
    \AnswerSpace{28mm}
    \item \(z_1=1+i\), \(z_2=4\left(\cos\frac{5\pi}{6}+i\sin\frac{5\pi}{6}\right)\). Určete \(\frac{z_2}{z_1}\) v algebraickém tvaru.
    \AnswerSpace{35mm}
    \item \(z=\sqrt2\left(\cos\frac{3\pi}{4}+i\sin\frac{3\pi}{4}\right)\). Pomocí Moivreovy věty vypočtěte \(z^4\) a výsledek zapište v algebraickém tvaru.
    \AnswerSpace{30mm}
\end{enumerate}

\newpage
\TestHeader{B}

% ZADÁNÍ VARIANTY B
\begin{enumerate}
    \item \((4-i)(2+5i)\)
    \AnswerSpace{17mm}
    \item \(\left|\frac{4-i}{2+2i}\right|\)
    \AnswerSpace{24mm}
    \item \(x^2+4x+5=0\), řešte v \(\mathbb C\).
    \AnswerSpace{28mm}
    \item \(z_1=1-i\), \(z_2=6\left(\cos\frac{2\pi}{3}+i\sin\frac{2\pi}{3}\right)\). Určete \(\frac{z_2}{z_1}\) v algebraickém tvaru.
    \AnswerSpace{35mm}
    \item \(z=2\left(\cos\frac{\pi}{6}+i\sin\frac{\pi}{6}\right)\). Pomocí Moivreovy věty vypočtěte \(z^3\) a výsledek zapište v algebraickém tvaru.
    \AnswerSpace{30mm}
\end{enumerate}
\end{document}
